\clearpage\thispagestyle{empty}
\pdfbookmark[0]{Quellen zur Besprechung von Klübers Mannheimer Sternwarte}{klueber1811-sources}
\begin{center}\textbf{Quellen zur Besprechung von Klübers Mannheimer Sternwarte}\end{center}
\noindent Diese KI-übersetzte Fassung in modernem Deutsch folgt Carl Friedrich Gauß, \textit{Werke}, Band~VI (Göttingen, 1874), S.~536--537. Die Besprechung erschien zuerst unter Mannheim und Heidelberg in den \textit{Göttingischen gelehrten Anzeigen}, Stück~192 und 193, 2. Dezember 1811, S.~1913--1916.
\par\smallskip\noindent In der Werkausgabe steht \textit{Astonomen} für \textit{Astronomen}. Vor dem Satz über das Bombardement fehlt der Punkt. Die Länge von Paris ist in der Werkausgabe mit hochgestelltem m und s, im Erstdruck mit einfachem und doppeltem Strich angegeben.
\par\smallskip\noindent Johann Ludwig Klübers \href{https://doi.org/10.11588/diglit.19120}{\textit{Die Sternwarte zu Mannheim}} (1811) beschreibt ein provisorisches Mittagsziel (\textit{Mire méridienne}) zum Ausrichten des Meridianfernrohrs. Die Stelle steht auf \href{https://doi.org/10.11588/diglit.19120#0061}{Scan~0061} des Exemplars der Universitätsbibliothek Heidelberg.
\par\smallskip\noindent \href{https://gdz.sub.uni-goettingen.de/id/PPN236007467}{Text der Werkausgabe}. Erstdruck, \href{https://images.p.sub.uni-goettingen.de/iiif/image/gdz:PPN319721507_1811:00001947/full/1650,/0/default.jpg}{S.~1913}, \href{https://images.p.sub.uni-goettingen.de/iiif/image/gdz:PPN319721507_1811:00001948/full/1650,/0/default.jpg}{S.~1914}, \href{https://images.p.sub.uni-goettingen.de/iiif/image/gdz:PPN319721507_1811:00001949/full/1650,/0/default.jpg}{S.~1915} und \href{https://images.p.sub.uni-goettingen.de/iiif/image/gdz:PPN319721507_1811:00001950/full/1650,/0/default.jpg}{S.~1916}.
\par\smallskip\noindent Übersetzung und Satz mit KI erstellt.
\clearpage
